\section{Compute Phase}

Once the model parameters are loaded, the design switches to the compute phase.
This phase again has two main pieces. The first is the input-preparation path
in \texttt{data\_in\_manager}, which converts the incoming activation stream
into the binary words expected by the network. The second is the layer-local
forward datapath inside \texttt{bnn\_layer}, where buffered activations, RAM
reads, and neuron-processor outputs are coordinated.

\subsection{Input Preparation}

\texttt{rtl/data\_in\_manager.sv} mirrors several of the same structural ideas
used in the configuration path. It begins with byte compaction on the external
stream, repacks the resulting bytes into bus-width words, applies the local
frame and padding rules needed by the design, binarizes pixel values, and then
stages the resulting binary words before handing them to the BNN core.

\begin{strip}
  \centering
  \resizebox{!}{0.39\textheight}{\begin{tikzpicture}[
  >=Latex,
  block/.style={
    draw,
    thick,
    rounded corners=1pt,
    minimum width=60mm,
    minimum height=17mm,
    align=center,
    font=\normalsize
  },
  wideblock/.style={
    block,
    minimum width=74mm
  },
  data/.style={->, thick},
  bus/.style={
    ->,
    thick,
    postaction={decorate},
    decoration={
      markings,
      mark=at position 0.55 with {
        \draw[thick] (-2.2pt,3pt) -- (2.2pt,-3pt);
      }
    }
  },
  ctrl/.style={->, thick, dashed},
  iface/.style={font=\large\bfseries\ttfamily, fill=white, inner sep=1pt},
  sig/.style={font=\normalsize\ttfamily, fill=white, inner sep=1pt},
  lab/.style={font=\normalsize, fill=white, inner sep=1pt}
]

\node[font=\Large\bfseries] at (0,3.3) {Data In Manager};

\node[block] (compact) at (0,0)
  {TKEEP Byte Compactor\\[-1pt]\texttt{tkeep\_byte\_compactor}};

\node[wideblock] (ctrlblk) at (0,-4.1)
  {Frame / Pad Control\\[-1pt]\textit{local logic}};

\node[block] (vw) at (0,-8.2)
  {VW Buffer\\[-1pt]\texttt{vw\_buffer}};

\node[block] (binarize) at (0,-12.4)
  {Binarizer\\[-1pt]\textit{byte $\geq$ 128}};

\node[block] (binfifo) at (0,-16.6)
  {Binary FIFO\\[-1pt]\texttt{bin\_fifo}};

% Top-level image inputs
\coordinate (divTop) at ([xshift=-48mm,yshift=18mm]compact.north);
\coordinate (divTurnA) at ([xshift=-48mm,yshift=8mm]compact.north);
\coordinate (divTurnB) at ([xshift=-20mm,yshift=8mm]compact.north);
\draw[data] (divTop) -- (divTurnA) -- (divTurnB) -- ([xshift=-20mm]compact.north);
\node[iface] at ($(divTop)+(0,2.5mm)$) {\texttt{data\_in\_valid}};

\draw[bus] ([yshift=18mm]compact.north) -- (compact.north);
\node[iface] at ([yshift=20.5mm]compact.north) {\texttt{data\_in\_data}};

\coordinate (dikTop) at ([xshift=48mm,yshift=18mm]compact.north);
\coordinate (dikTurnA) at ([xshift=48mm,yshift=8mm]compact.north);
\coordinate (dikTurnB) at ([xshift=20mm,yshift=8mm]compact.north);
\draw[bus] (dikTop) -- (dikTurnA) -- (dikTurnB) -- ([xshift=20mm]compact.north);
\node[iface] at ($(dikTop)+(0,2.5mm)$) {\texttt{data\_in\_keep}};

% Compactor -> frame/pad control
\draw[data] ([xshift=-22mm]compact.south) -- ([xshift=-22mm]ctrlblk.north)
  node[pos=0.64,left,sig] {\texttt{compact\_wr\_en}};

\draw[bus] (compact.south) -- (ctrlblk.north)
  node[pos=0.40,right,sig,xshift=2pt] {\texttt{compact\_wr\_data}};

\draw[bus] ([xshift=22mm]compact.south) -- ([xshift=22mm]ctrlblk.north)
  node[pos=0.64,right,sig] {\texttt{compact\_total\_bytes}};

% Side inputs into frame/pad control
\draw[data] ([xshift=-38mm,yshift=7mm]ctrlblk.west) -- ([yshift=7mm]ctrlblk.west)
  node[pos=0.48,above,iface] {\texttt{config\_ready}};

\draw[data] ([xshift=-38mm]ctrlblk.west) -- (ctrlblk.west)
  node[pos=0.48,above,iface] {\texttt{config\_last}};

\draw[data] ([xshift=-38mm,yshift=-7mm]ctrlblk.west) -- ([yshift=-7mm]ctrlblk.west)
  node[pos=0.48,below,iface] {\texttt{data\_in\_last}};

% data_in_ready routed back up as interface signal
\coordinate (dinreadyA) at ($(ctrlblk.east)+(4.0,0)$);
\coordinate (dinreadyTop) at ($(dinreadyA |- compact.north)+(0,18mm)$);

\draw[thick] (ctrlblk.east) -- (dinreadyA);
\draw[data] (dinreadyA) -- (dinreadyTop);
\node[iface] at ($(dinreadyTop)+(0,2.3mm)$) {\texttt{data\_in\_ready}};

% Frame/pad control -> vw_buffer
\draw[data] ([xshift=-22mm]ctrlblk.south) -- ([xshift=-22mm]vw.north)
  node[pos=0.64,left,sig] {\texttt{vw\_issue\_en}};

\draw[bus] (ctrlblk.south) -- (vw.north)
  node[pos=0.40,right,sig,xshift=2pt] {\texttt{vw\_issue\_data}};

\draw[bus] ([xshift=22mm]ctrlblk.south) -- ([xshift=22mm]vw.north)
  node[pos=0.64,right,sig] {\texttt{vw\_issue\_total\_bytes}};

% vw_buffer -> binarizer / fifo
\draw[bus] (vw.south) -- (binarize.north)
  node[pos=0.58,right,sig] {\texttt{vw\_rd\_data}};

\draw[data] ([xshift=-34mm]vw.south) -- ([xshift=-34mm]binfifo.north)
  node[pos=0.34,left,sig] {\texttt{vw\_rd\_en}};

% Binarizer -> binary FIFO
\draw[bus] (binarize.south) -- (binfifo.north)
  node[pos=0.56,right,sig] {\texttt{bin\_fifo\_wr\_data}};

% Binary FIFO feedback to control
\coordinate (almfullA) at ($(binfifo.east)+(3.8,0)$);
\coordinate (almfullB) at ($(almfullA |- ctrlblk.east)$);

\draw[ctrl] (binfifo.east) -- (almfullA) -- (almfullB) -- (ctrlblk.east);
\node[sig,right] at ($(almfullA)!0.55!(almfullB)$) {\texttt{bin\_fifo\_alm\_full}};

% BNN-side control into fifo staging
\draw[data] ([xshift=38mm,yshift=6mm]binfifo.east) -- ([yshift=6mm]binfifo.east)
  node[pos=0.48,above,iface] {\texttt{bnn\_ready}};

\draw[data] ([xshift=38mm,yshift=-6mm]binfifo.east) -- ([yshift=-6mm]binfifo.east)
  node[pos=0.48,below,iface] {\texttt{bnn\_en}};

% Outputs to BNN
\draw[bus] ([xshift=-17mm]binfifo.south) -- ++(0,-2.6)
  node[below,iface] {\texttt{bnn\_data\_in}};

\draw[data] ([xshift=17mm]binfifo.south) -- ++(0,-2.6)
  node[below,iface] {\texttt{bnn\_data\_in\_valid}};

\end{tikzpicture}
}
  \captionof{figure}{Block diagram of the \texttt{data\_in\_manager} datapath
  and its primary interface signals.}
  \Description{Block diagram of the data input manager showing valid-byte
  compaction, repacking, input framing, binarization, and the staged interface
  that feeds the BNN core.}
  \label{fig:data-in-manager}
\end{strip}

This preprocessing stage is important because it decouples the external input
stream shape from the internal datapath width. By the time activations leave
\texttt{data\_in\_manager}, they already match the first layer's expected
binary format and can be replayed locally without repeatedly revisiting the
original byte-level protocol.

\clearpage

\subsection{Layer Compute Flow}

Inside \texttt{rtl/bnn\_layer.sv}, the forward path is organized around the
replay buffer, the read-side control logic, the banked parameter memories, and
the neuron-processor array. The intent is to keep the activation stream and the
shared RAM-read schedule locally synchronized while still exposing the timing
registers that matter to the implementation.

\begin{figure*}[t]
  \centering
  \resizebox{0.95\textwidth}{!}{\begin{tikzpicture}[
  >=Latex,
  block/.style={
    draw,
    thick,
    rounded corners=1pt,
    minimum width=74mm,
    minimum height=26mm,
    align=center,
    font=\large
  },
  tallblock/.style={
    block,
    minimum width=72mm,
    minimum height=60mm
  },
  smallblock/.style={
    draw,
    thick,
    rounded corners=1pt,
    minimum width=42mm,
    minimum height=20mm,
    align=center,
    font=\large
  },
  wideblock/.style={
    block,
    minimum width=98mm,
    minimum height=32mm
  },
  data/.style={->, thick},
  bus/.style={
    ->,
    thick,
    postaction={decorate},
    decoration={
      markings,
      mark=at position 0.55 with {
        \draw[thick] (-2.2pt,3pt) -- (2.2pt,-3pt);
      }
    }
  },
  ctrl/.style={->, thick, dashed},
  iface/.style={font=\large\bfseries\ttfamily, fill=white, inner sep=1pt},
  sig/.style={font=\normalsize\ttfamily, fill=white, inner sep=1pt},
  portsig/.style={font=\normalsize\bfseries\ttfamily, fill=white, inner sep=1pt, anchor=east, align=right}
]

\node[font=\Large\bfseries] at (2.8,5.7) {BNN Layer Compute Flow};

\node[wideblock] (cfg) at (0,0)
  {Config Path\\[-1pt]\textit{staging + addr gen}};

\node[smallblock] (wwrite) at (7.0,-4.5)
  {W Write\\Reg};

\node[smallblock] (twrite) at (13.0,-4.5)
  {T Write\\Reg};

\node[block] (replay) at (-14,-9.6)
  {Replay Buffer\\[-1pt]\texttt{replay\_buffer}};

\node[smallblock] (align) at (-6.4,-9.6)
  {Align\\Delay};

\node[block] (wram) at (7.0,-9.6)
  {Weight RAM Banks\\[-1pt]\texttt{u\_w\_ram[gi]}};

\node[wideblock] (nc) at (-14,-15.0)
  {Neuron Controller\\[-1pt]\texttt{neuron\_controller}};

\node[smallblock] (wread) at (-0.8,-9.6)
  {W Read\\Regs};

\node[smallblock] (tread) at (-0.8,-15.0)
  {T Read\\Regs};

\node[block] (tram) at (7.0,-15.0)
  {Threshold RAM Banks\\[-1pt]\texttt{u\_t\_ram[gi]}};

\node[tallblock] (np) at (21.0,-12.0)
  {Neuron Processors\\[-1pt]\texttt{u\_np[gi]}\\[-1pt]\textit{+ output regs}};

\coordinate (npXin) at ([yshift=7mm]np.west);
\coordinate (npWin) at ([yshift=3mm]np.west);
\coordinate (npTin) at ([yshift=-5mm]np.west);
\coordinate (npVin) at ([yshift=-10mm]np.west);
\coordinate (npLastin) at ([yshift=-14mm]np.west);
\coordinate (xEntry) at ($(npXin)+(-4mm,0)$);
\coordinate (wEntry) at ($(npWin)+(-12mm,0)$);
\coordinate (tEntry) at ($(npTin)+(-10mm,0)$);
\coordinate (vEntry) at ($(npVin)+(-8mm,0)$);
\coordinate (lastEntry) at ($(npLastin)+(-6mm,0)$);

% Layer interface inputs
\draw[bus] ([xshift=-30mm,yshift=18mm]cfg.north) -- ([xshift=-30mm]cfg.north);
\node[iface] at ([xshift=-30mm,yshift=20.5mm]cfg.north) {\texttt{w\_data}};

\draw[data] ([xshift=-10mm,yshift=18mm]cfg.north) -- ([xshift=-10mm]cfg.north);
\node[iface] at ([xshift=-10mm,yshift=20.5mm]cfg.north) {\texttt{w\_en}};

\draw[bus] ([xshift=10mm,yshift=18mm]cfg.north) -- ([xshift=10mm]cfg.north);
\node[iface] at ([xshift=10mm,yshift=20.5mm]cfg.north) {\texttt{t\_data}};

\draw[data] ([xshift=30mm,yshift=18mm]cfg.north) -- ([xshift=30mm]cfg.north);
\node[iface] at ([xshift=30mm,yshift=20.5mm]cfg.north) {\texttt{t\_en}};

\draw[bus] ([yshift=4mm]replay.west) -- ++(-4.0,0);
\node[iface] at ($(replay.west)+(-23mm,4mm)$) {\texttt{data\_in}};

\draw[data] ([yshift=-4mm]replay.west) -- ++(-4.0,0);
\node[iface] at ($(replay.west)+(-23mm,-4mm)$) {\texttt{valid\_in}};

\coordinate (rdyLane) at ($(replay.north west)+(-1.8,0)$);
\coordinate (rdyTop) at ($(rdyLane |- cfg.north)+(0,10mm)$);
\draw[thick] (replay.north west) -- (rdyLane);
\draw[data] (rdyLane) -- (rdyTop);
\node[iface] at ($(rdyTop)+(0,2.3mm)$) {\texttt{ready\_in}};

% Layer outputs
\draw[data] ([yshift=10mm]np.east) -- ++(4.6,0);
\node[iface] at ($(np.east)+(24mm,10mm)$) {\texttt{valid\_out}};

\draw[bus] (np.east) -- ++(4.6,0);
\node[iface] at ($(np.east)+(24mm,0mm)$) {\texttt{data\_out}};

\draw[bus] ([xshift=10mm]np.south) -- ++(0,-2.8);
\node[iface] at ([xshift=10mm,yshift=-3.4mm]np.south) {\texttt{count\_out}};

% Config write path to RAMs
\coordinate (cfgW0) at ([xshift=14mm]cfg.south);
\coordinate (cfgW1) at ($(wwrite.north)+(0,0.9)$);
\coordinate (cfgW2) at (cfgW1 -| cfgW0);
\draw[ctrl] (cfgW0) -- (cfgW2) -- (cfgW1) -- (wwrite.north);

\coordinate (cfgT0) at ([xshift=28mm]cfg.south);
\coordinate (cfgT1) at ($(twrite.north)+(0,0.9)$);
\coordinate (cfgT2) at (cfgT1 -| cfgT0);
\draw[ctrl] (cfgT0) -- (cfgT2) -- (cfgT1) -- (twrite.north);

\draw[bus] (wwrite.south) -- (wram.north);

\coordinate (tWriteIn) at ([xshift=18mm]tram.north);
\coordinate (tWriteA) at ($(tWriteIn)+(0,1.1)$);
\coordinate (tWriteB) at (tWriteA -| twrite.south);
\draw[bus] (twrite.south) -- (tWriteB) -- (tWriteA) -- (tWriteIn);

% Replay path
\draw[bus] (replay.east) -- (align.west)
  node[pos=0.50,above,sig] {\texttt{x\_buf}};

\coordinate (xLaneA) at ([xshift=2mm]align.east);
\coordinate (xLaneB) at ($(xLaneA)+(0,28mm)$);
\coordinate (xLaneCBase) at ([xshift=-10mm]wwrite.west);
\coordinate (xLaneC) at (xLaneCBase |- xLaneB);
\coordinate (xLaneD) at ($(xLaneC |- wwrite.north)+(0,5mm)$);
\coordinate (xLaneE) at (xLaneD -| xEntry);
\draw[bus] (align.east) -- (xLaneA) -- (xLaneB) -- (xLaneC) -- (xLaneD) -- (xLaneE) -- (xEntry) -- (npXin);

% Replay / config control into neuron controller
\draw[data] (replay.south) -- (nc.north)
  node[pos=0.52,left,sig] {\texttt{issue}};

\coordinate (doneA) at ($(cfg.south west)+(-4.5,0)$);
\coordinate (doneB) at (doneA |- nc.north);
\coordinate (doneC) at ([xshift=-22mm]nc.north);
\draw[ctrl] (cfg.south west) -- (doneA) -- (doneB) -- (doneC) -- (nc.north);
\node[sig,left] at ($(doneA)!0.50!(doneB)$) {\texttt{cfg\_done}};

% Read address path
\coordinate (wReadSide) at ([yshift=4mm]wread.west);
\coordinate (tReadSide) at ([yshift=-4mm]tread.west);
\coordinate (addrLaneTop) at ([xshift=-9mm,yshift=4mm]wread.west);
\coordinate (addrLaneBot) at ([xshift=-9mm,yshift=-4mm]tread.west);

\coordinate (wAddr0) at ([yshift=4mm]nc.east);
\coordinate (wAddr1) at (addrLaneTop |- wAddr0);
\draw[data] (wAddr0) -- (wAddr1) -- (addrLaneTop) -- (wReadSide);
\node[sig,above] at ($(wAddr0)!0.55!(wAddr1)$) {\texttt{w\_addr}};

\coordinate (tAddr0) at ([yshift=-4mm]nc.east);
\coordinate (tAddr1) at (addrLaneBot |- tAddr0);
\draw[data] (tAddr0) -- (tAddr1) -- (addrLaneBot) -- (tReadSide);
\node[sig,below] at ($(tAddr0)!0.55!(tAddr1)$) {\texttt{t\_addr}};

\draw[data] (wread.east) -- (wram.west);

\draw[data] (tread.east) -- (tram.west)
  node[pos=0.50,below,sig] {\texttt{t regs}};

% RAM data to neuron processors
\coordinate (wToNpStart) at ([yshift=4mm]wram.east);
\coordinate (wToNpA) at (wToNpStart -| wEntry);
\draw[bus] (wToNpStart) -- (wToNpA) -- (wEntry) -- (npWin);

\coordinate (tToNpA) at (tram.east -| tEntry);
\draw[bus] (tram.east) -- (tToNpA) -- (tEntry) -- (npTin);

% Controller -> neuron processors
\coordinate (vLane0) at ([yshift=-6mm]nc.east);
\coordinate (vLane1) at ($(vLane0)+(3.2,0)$);
\coordinate (vLane2) at (vLane1 |- vEntry);
\draw[data] (vLane0) -- (vLane1) -- (vLane2) -- (vEntry) -- (npVin);

\coordinate (last0) at ([yshift=-12mm]nc.east);
\coordinate (last1) at ($(last0)+(3.2,0)$);
\coordinate (last2) at (last1 |- lastEntry);
\draw[data] (last0) -- (last1) -- (last2) -- (lastEntry) -- (npLastin);

% Processor-side interface labels
\node[portsig] at ($(npXin)+(-2mm,0)$) {\texttt{x}};
\node[portsig] at ($(npWin)+(-2mm,0)$) {\texttt{weights}};
\node[portsig] at ($(npTin)+(-2mm,0)$) {\texttt{thresh}};
\node[portsig] at ($(npVin)+(-2mm,0)$) {\texttt{v\_in}};
\node[portsig] at ($(npLastin)+(-2mm,0)$) {\texttt{last}};

\end{tikzpicture}
}
  \caption{Simplified compute flow inside one \texttt{bnn\_layer}. The replay
  buffer, read-control path, RAM banks, and neuron-processor array are shown in
  the order they participate in inference.}
  \Description{Block diagram of the layer compute flow showing replay buffering
  of activations, alignment, controller-driven reads from weight and threshold
  RAM banks, and the final neuron-processor array that produces outputs.}
  \label{fig:bnn-layer-compute}
\end{figure*}

The replay buffer and align-delay logic keep the activation side available when
the local controller issues parameter reads. In parallel,
\texttt{neuron\_controller} sequences the weight and threshold-bank access
pattern, and the neuron processors combine those operands to generate the next
layer's outputs. Together, these blocks form the runtime half of the design
that follows the earlier configuration phase.
